\ficha{Emissão e suspensão do troco, 1914}

\bloco{Fontes lidas}

Não há volume parlamentar para 1914. Li, na imprensa: o resumo da sessão da
Câmara no \textit{Correio Paulistano} de 31 de julho; o \textit{Correio da
Manhã} de 5 de agosto, com a entrevista de Leopoldo de Bulhões e o discurso de
Rui Barbosa no Senado; \textit{O Paiz} de 6, 11, 12, 18 e 21 de agosto; e a
\textit{Gazeta de Notícias} de 12 de dezembro, com a discussão do projeto que
prorrogou até 31 de dezembro de 1915 a suspensão do troco decretada em 15 de
agosto \cite{brasil1914suspensao}.

\bloco{Síntese}

A guerra europeia encontra o Tesouro em crise. Em agosto o governo decreta
feriado bancário, o Congresso vota a moratória, a emissão de papel-moeda é
autorizada \cite{brasil1914emissao} e a Caixa deixa de trocar suas notas à
vista. Há dois
objetos de desacordo, a emissão e o troco, e as posições sobre um não predizem
as posições sobre o outro. Contra a suspensão do troco aparecem duas razões, a
honra do depositário e o desejo de ver a Caixa acabar, que não separam pessoas:
Serzedelo Corrêa usa as duas na mesma fala.

\bloco{O troco: a Caixa como depositária}

Em entrevista ao \textit{Correio da Manhã}, Bulhões atribui a baixa do câmbio
às dificuldades postas à retirada do ouro e formula o princípio. Bancos
emissores, com lastro parcial, podem se defender numa crise; a Caixa, não:

\cit[I:per089842_1914_05641:2]{Ella deve ser fiel depositaria. As suas notas
correspondem a titulos de deposito}{\jor{cm19140805}{Correio da Manhã}{5 ago.
1914}{2}}

Dias antes, segundo o resumo do \textit{Correio Paulistano}, Calógeras dizia
que o governo não podia lançar mão de medidas arbitrárias como
\tr[I:per090972_1914_18331:4]{impedir as retiradas da Caixa de Conversão}, e
que a emissão seria \tr[I:per090972_1914_18331:4]{um verdadeiro crime}
(\jor{cp19140731}{Correio Paulistano}{31 jul. 1914}{4}).

Em dezembro, na Câmara, Martim Francisco recusa o projeto do Senado porque ele
autoriza o governo, como depositário, a não restituir o que lhe foi confiado.
Seria aconselhar a \tr[I:per103730_1914_00345:2]{pátria brasileira a ser
deshonesta} (\jor{gn19141212}{Gazeta de Notícias}{12 dez. 1914}{2}). Serzedelo
Corrêa o acompanha e, na mesma intervenção, diz duas coisas:

\cit[I:per103730_1914_00345:2]{Suspender a retirada de ouro é
deshonesto}{\jor{gn19141212}{Gazeta de Notícias}{12 dez. 1914}{2}}

\cit[I:per103730_1914_00345:2]{A Caixa é um grande trambolho}{\jor{gn19141212}{Gazeta
de Notícias}{12 dez. 1914}{2}}

Conferi na imagem que o trecho está sob o título com o nome de Serzedelo, em
coluna própria. Ele quer que a Caixa funcione normalmente e, se os depósitos se
esgotarem e ela desaparecer, tanto melhor.

\bloco{A emissão}

Serzedelo se declara contrário à emissão por doutrina e a aceita por
necessidade:

\cit[I:per178691_1914_10901:4]{o actual governo não tem hoje outra saida senão a
emissão do papel-moeda}{\jor{opaiz19140812}{O Paiz}{12 ago. 1914}{4}}

Acrescenta que a emissão \tr[I:per178691_1914_10901:4]{terá a sua acção malefica
um pouco minorada} pela existência da Caixa, cujas notas os bancos haviam
convertido em ouro. Martim Francisco vota \tr[I:per178691_1914_10900:4]{contra a
emissão de papel-moeda} (\jor{opaiz19140811}{O Paiz}{11 ago. 1914}{4}).
Cincinato Braga, também paulista, justifica a atitude de São Paulo,
\tr[I:per178691_1914_10910:6]{favoravel á emissão de papel-moeda}
(\jor{opaiz19140821}{O Paiz}{21 ago. 1914}{6}).

Antônio Carlos, relator vencido na Comissão de Finanças, combate o projeto do
Senado e propõe letras do Tesouro. \textit{O Paiz} publica seu voto e, na mesma
página, um texto da redação que o rebate. O jornal resume a acusação do
deputado aos bancos estrangeiros, que teriam enfraquecido as próprias caixas

\cit[I:per178691_1914_10907:3]{trocando as suas notas pelo ouro da Caixa de
Conversão, destinado á remessa para as suas matrizes}{\jor{opaiz19140818}{O
Paiz}{18 ago. 1914}{3}}

e defende a emissão: ela \tr[I:per178691_1914_10907:3]{foi o unico alvitre
pratico que surgiu para fazer face ás difficuldades de momento}.

\bloco{Rui Barbosa}

A fala direta de Rui que localizei foi publicada em 5 de agosto e trata do decreto do
feriado, não da emissão. Ele o considera ilegal por invadir atribuição do
Congresso:

\cit[I:per089842_1914_05641:2]{o perigo monstruoso da usurpação desta
competencia legislativa}{\jor{cm19140805}{Correio da Manhã}{5 ago. 1914}{2}}

\textit{O Paiz}, em texto da redação, o censura por não propor saída para a
crise (\jor{opaiz19140806}{O Paiz}{6 ago. 1914}{3}). Sobre a emissão, a posição
de Rui só aparece de modo indireto, em resumo da fala de Bulhões no Senado
(\jor{cm19140808}{Correio da Manhã}{8 ago. 1914}{2}).

\bloco{Achados}

\begin{enumerate}[leftmargin=*,nosep]
\item O argumento do depositário tem em 1914 a formulação direta de Bulhões e
continua o que ele sustentava em 1910 sobre a propriedade do ouro (ficha 5) e
o que \textit{O Paiz} escrevia em 1909 (ficha 4).
\item As duas razões contra a suspensão do troco coexistem na fala de
Serzedelo. O que o distingue de Martim Francisco é desejar o fim da Caixa, não
a ausência do argumento da honra.
\item São Paulo se divide sobre a emissão: Cincinato Braga a favor, Martim
Francisco contra.
\item A objeção de Rui Barbosa em agosto é de competência entre os poderes.
Não há, na base, fala dele sobre a emissão.
\end{enumerate}

\bloco{Limites}

Todas as falas de 1914 vêm de jornal. Só a entrevista de Bulhões e o discurso
de Rui estão em primeira pessoa; as demais são resumos de sessão. A data da
edição do \textit{Correio Paulistano} com a fala de Calógeras foi lida no OCR
do cabeçalho e não conferida na imagem.
